Ce projet vise à utiliser de grands modèles de langage (LLM) pour l’analyse vidéo.

À ce jour, le traitement automatique de longs segments vidéo requiert souvent l'utilisation de modèles multimodaux complexes, qui exigent une puissance de calcul élevée, peu compatible avec les capacités matérielles limitées des robots mobiles.

Ce projet  consiste à représenter intégralement les informations vidéo sous la forme d'une description textuelle détaillée et structurée en reposant exclusivement sur une représentation textuelle exhaustive, englobant les dialogues, les événements visuels, l'ambiance sonore et même les expressions émotionnelles. L'objectif est de démontrer que cette méthode pourrait simplifier considérablement l'analyse vidéo, en réduisant la dépendance aux modèles visuels et multi-modaux habituellement utilisés. La performance du modèle développé serait évaluée en comparant systématiquement les moments saillants trouvés par le modèle à des segments de référence sélectionnés par des humains, à partir d'une base de données de vidéos Youtube et de leur extraits populaires correspondants. Le choix de cette base de données se justifie par l'abondance de données disponibles et leur pertinence pour la tâche visée qui est l’analyse de vidéos de tous genre en robotique.

La principale contribution scientifique de la recherche est la création d’un standard de représentation textuelle de vidéos qui permetterait l’analyse complète du vidéo par un grand modèle de langage textuel, permettant la détection automatisée des événements importants d’un vidéo. En évitant le recours à des méthodes d'analyse multimodales intensives, cette approche pourrait permettre l'intégration d'une détection efficace des moments forts sur des dispositifs robotiques soumis à des contraintes computationnelles strictes. La recherche apporterait donc une contribution directe au monde de la robotique facilitant l'exploitation des grands modèles de langage pour le traitement de données vidéo complexes dans des contextes embarqués.
